\question{}

\begin{ابجد}
\item \textbf{توجه:} بر اساس الگوریتم \code{TREE-SEARCH} موجود در اسلایدهای درس و کتاب راسل و نورویگ فرض می‌کنیم هر گرهی که از \lr{fringe} خارج می‌شود، اگر هدف نباشد بسط داده می‌شود، حتی اگر خودش برگ باشد و فرزندی نداشته باشد. \\
چون درخت کامل است، می‌دانیم که تعداد گره‌ها در عمق دلخواه $x$ برابر با $b^x$ بوده و تعداد کل گره‌های درخت نیز از $O(b^d)$ است. حال هر الگوریتم را جداگانه بررسی می‌کنیم:
\begin{itemize}
    \item الگوریتم \lr{DFS}:
          \begin{description}
              \item[بهترین حالت:] می‌دانیم که \lr{DFS} یک مسیر را انتخاب کرده و آن را به صورت عمقی تا جایی ادامه می‌دهد که به هدف یا بن‌بست بخورد. در بهترین حالت، این الگوریتم در همان ابتدا مسیر درست را انتخاب کرده و به هدف می‌رسد. بنابراین {$m$ گره} را بسط می‌دهد. (خود هدف به عنوان گره $(m+1)$ ام از \lr{fringe} خارج می‌شود و چون هدف است دیگر بسط داده نمی‌شود.) \\
                  این جواب از $O(m)$ است.
              \item[بدترین حالت:] در بدترین حالت، \lr{DFS} تمام مسیرهای دیگر را تا انتها (عمق $d$) پیش می‌رود و در انتها مسیر درست را انتخاب می‌کند. پس در این حالت، الگوریتم تمام گره‌های درخت به جز خود گره هدف و زیردرخت آن را بسط می‌دهد. در این صورت تعداد کل گره‌های بسط داده شده برابر است با:
                  \begin{align*}
                      \frac{b^{d+1} - 1}{b-1} - \frac{b^{d-m+1} - 1}{b-1} = \frac{b^{d+1} - b^{d-m+1}}{b-1}
                  \end{align*}
                  که جواب از {$O(b^d)$} می‌باشد.
          \end{description}

    \item الگوریتم \lr{BFS}:
          \begin{description}
              \item[بهترین حالت:] \lr{BFS} همواره کم عمق‌ترین عضو را از \lr{fringe} خارج می‌کند؛ در واقع می‌توانیم بگوییم \lr{BFS} به صورت لایه‌ای پیش می‌رود: یعنی ابتدا تمام گره‌ها با عمق 1 را بررسی می‌کند، سپس تمام گره‌ها با عمق 2 و \dots \\
                  بنابراین در بهترین حالت، اولین گره با عمق $m$ که از \lr{fringe} خارج می‌شود همان هدف است. در این حالت، تعداد گره‌های بسط داده شده برابر با یک درخت کامل با ضریب انشعابی $b$ و عمق $m-1$ است که تعداد گره‌های آن با تصاعد هندسی محاسبه شده و برابر با $\frac{b^m - 1}{b - 1}$ است. \\
                  جواب از $O(b^{m-1})$ است.
              \item[بدترین حالت:] در این حالت، هدف آخرین گره از لایه با عمق $m$ است که توسط \lr{BFS} از \lr{fringe} خارج می‌شود. بنابراین علاوه بر تعدادی که در مورد قبل محاسبه کردیم، $b^m - 1$ عضو دیگر نیز بسط داده می‌شود و تعداد کل گره‌های بسط داده شده در این حالت برابر است با:
                  \begin{align*}
                      \frac{b^m - 1}{b - 1} + b^m - 1 = \frac{b^{m+1} - b}{b - 1}
                  \end{align*}
                  که جواب از $O(b^m)$ می‌باشد.
          \end{description}
\end{itemize}

\item فرض می‌کنیم در هر دو الگوریتم \lr{BFS} و \lr{DFS} هنگامی که به گره‌هایی با اولویت یکسان می‌رسند، گره چپ‌تر را زودتر انتخاب می‌کنند.
\begin{itemize}
    \item \textbf{چپ‌ترین گره}:
          \begin{itemize}
              \item[] این همان بهترین حالت \lr{DFS} است که در بخش قبل محاسبه کردیم؛ الگوریتم از ریشه درخت شروع کرده و مستقیم به سمت گره هدف می‌رود. در کل $m+1$ گره از \lr{fringe} خارج می‌شوند.
              \item[] \lr{BFS} تمام گره‌های لایه‌های بالاتر را بررسی می‌کند تا به گره هدف برسد. یعنی در کل $\frac{b^m - 1}{b - 1} + 1$ گره از \lr{fringe} خارج می‌شود.
              \item[] در این حالت، عملکرد الگوریتم \lr{DFS} به مراتب بهتر است.
          \end{itemize}

    \item \textbf{راست‌ترین گره}:
          \begin{itemize}
              \item[] این روش بدترین حالت \lr{DFS} است و مطابق بخش (آ) تعداد گره‌های خارج شده از $O(b^d)$ می‌باشد.
              \item[] در این حالت گره هدف آخرین گره از لایه به عمق $m$ است که از \lr{fringe} خارج می‌شود. در این حالت تعداد کل گره‌های خارج شده، برابر $\frac{b^{m+1}-1}{b-1}$ است که از $O(b^m)$ می‌باشد.
              \item[] در این حالت در صورتی که هدف در لایه آخر نباشد ($m<d$) الگوریتم \lr{BFS} عملکرد بهتری دارد.
          \end{itemize}

    \item \textbf{حالت متوسط}:
          \begin{itemize}
              \item[] می‌خواهیم امید ریاضی تعداد گره‌های خارج شده از \lr{fringe} (شامل خود هدف) تا زمان رسیدن به هدف را محاسبه کنیم. برای سادگی، می‌توانیم احتمال خارج شدن هر گره را محاسبه کرده و با هم جمع بزنیم. برای این کار، گره‌ها را بر اساس عمق آنها ($h$) به دو دسته تقسیم می‌کنیم:
                  \begin{itemize}
                      \item[] $h \leq m$: در هر یک از این لایه‌ها $b^h$ گره داریم. این گره‌ها در صورتی خارج می‌شوند که یا یکی از اجداد گره هدف باشند (چون در هر لایه بالاتر از هدف، دقیقا یکی از گره‌ها جد گره هدف است. پس احتمال اینکه یک گره تصادفی در لایه $h$ جد هدف باشد برابر با $\frac{1}{b^h}$ است.) و یا سمت چپ جد گره هدف بوده‌اند و زودتر از آن از \lr{fringe} خارج شده‌اند(به احتمال $\frac{1}{2}(1 - \frac{1}{b^h})$). اگر این احتمال را برای تمام این گره‌ها جمع بزنیم خواهیم داشت:
                          \begin{align*}
                              \sum_{h=0}^{m} b^h(\frac{1}{2} + \frac{1}{2b^h}) = \frac{m+1}{2} + \frac{b^{m+1}-1}{2(b-1)}
                          \end{align*}
                      \item[] $h > m$: هر گرهی که در این لایه‌ها قرار دارد، یک جد در عمق $m$ دارد، که این جد به احتمال برابر یکی از $b^m$ گره عمق مورد نظر است. این گره در صورتی پیش از هدف از \lr{fringe} خارج می‌شود که جد آن در عمق $m$ سمت چپ هدف قرار داشته باشد (به احتمال $\frac{1}{2}(1 - \frac{1}{b^m})$). این احتمال‌ها را جمع می‌زنیم:
                          \begin{align*}
                              \sum_{h=m+1}^{d} b^h(\frac{1}{2})(1 - \frac{1}{b^m}) = \frac{b(b^m-1)(b^{d-m}-1)}{2(b-1)}
                          \end{align*}
                  \end{itemize}
                  این دو مقدار را از نظر اردر جمع می‌زنیم:
                  \begin{align*}
                      O(\frac{m+1}{2} + \frac{b^{m+1}-1}{2(b-1)}) + O(\frac{b(b^m-1)(b^{d-m}-1)}{2(b-1)}) & = O(m + b^m + b^d) \\
                      \xRightarrow{m \leq d}                                                              & = O(b^d)
                  \end{align*}
              \item[] می‌دانیم که \lr{BFS} ابتدا تمام گره‌های لایه‌های بالاتر را از \lr{fringe} خارج می‌کند ($\frac{b^m - 1}{b - 1}$ عضو). در لایه به عمق $m$، $b^m$ گره داریم. فرض می‌کنیم گره هدف با احتمال برابر یکی از این گره‌ها باشد، امید ریاضی تعداد گره‌هایی که از لایه آخر خارج می‌شوند تا آخرینِ آنها گره هدف باشد را محاسبه می‌کنیم:
                  \begin{align*}
                      \frac{1}{b^m} + \frac{2}{b^m} + \dots + \frac{b^m}{b^m} & = \frac{1}{b^m}(1 + 2 + \dots + b^m)                        \\
                                                                              & = \frac{1}{b^m}(\frac{b^m(b^m + 1)}{2}) = \frac{b^m + 1}{2}
                  \end{align*}
                  که از $O(b^m)$ است. تعداد گره‌های خارج شده در لایه‌های بالاتر هم از $O(b^{m-1})$ بود؛ پس در کل $O(b^m)$ گره از \lr{fringe} خارج می‌شوند تا به گره هدف برسیم.
              \item[] اگر هدف در لایه آخر نباشد ($m<d$)، الگوریتم \lr{BFS} به طور متوسط عملکرد بهتری نسبت به \lr{DFS} دارد. البته بر اساس قسمت‌های قبل می‌توان به طور شهودی گفت که اگر هدف سمت چپ درخت باشد، \lr{DFS} عملکرد بهتری دارد و هر چه هدف به سمت راست حرکت کند، عملکرد \lr{DFS} کاهش می‌یابد تا زمانی که عملکرد \lr{BFS} بهتر از آن شود.
          \end{itemize}
\end{itemize}

\item بر اساس امید ریاضی‌هایی که در قسمت قبل محاسبه کردیم، در این مثال تا زمانی که هدف در لایه آخر نباشد، عملکرد \lr{BFS} بهتر است؛ اگر هدف در لایه آخر باشد، تعداد گره‌های بسط داده شده در هر دو الگوریتم از یک مرتبه است.

\end{ابجد}